\documentclass[14pt,a4paper]{extarticle}

\usepackage{fontspec}
\usepackage{polyglossia}
\setdefaultlanguage{russian}
\setotherlanguage{english}

\setmainfont{Tinos}
\setsansfont{Arimo}
\setmonofont{Cousine}[Scale=0.82]

\usepackage{geometry}
\geometry{a4paper, top=2cm, bottom=2cm, left=3cm, right=1.5cm}

\usepackage{amsmath}
\usepackage{amssymb}
\usepackage{graphicx}
\usepackage{float}
\usepackage{booktabs}
\usepackage{xcolor}
\usepackage{listings}
\usepackage{caption} 
\usepackage{tikz}
\usetikzlibrary{positioning}
\usepackage{enumitem}
\usepackage{hyperref}

\hypersetup{hidelinks, unicode}

\linespread{1.3}
\emergencystretch=1.5em
\setlength{\parindent}{1.25cm}

\captionsetup[figure]{name=Рисунок, labelsep=endash}
\captionsetup[table]{name=Таблица, labelsep=endash, singlelinecheck=false, justification=raggedright}
\renewcommand{\lstlistingname}{Листинг}

\newcommand{\Student}{А. А. Устинов}
\newcommand{\Group}{М8О-307БВ-24}
\newcommand{\Teacher}{В. Д. Бахарев}
\newcommand{\Year}{2026}

\definecolor{codegray}{gray}{0.45}
\definecolor{codegreen}{rgb}{0.0,0.45,0.2}
\definecolor{codeblue}{rgb}{0.1,0.2,0.65}

\lstdefinestyle{code}{
	basicstyle=\ttfamily\footnotesize,
	keywordstyle=\color{codeblue}\bfseries,
	commentstyle=\color{codegreen}\itshape,
	stringstyle=\color{codegray},
	numbers=left,
	numberstyle=\tiny\color{codegray},
	numbersep=8pt,
	breaklines=true,
	breakatwhitespace=false,
	showstringspaces=false,
	tabsize=4,
	extendedchars=true,
	inputencoding=utf8,
	frame=single,
	rulecolor=\color{codegray},
	framesep=4pt,
	captionpos=b,
	literate={\ }{{\ }}1
}
\lstset{style=code, language=C++}

\lstdefinelanguage{GLSL}[]{C++}{
	morekeywords={vec2,vec3,vec4,mat2,mat3,mat4,layout,location,uniform,in,out,
		push_constant,gl_Position,mix,version},
}

\newcommand{\screenshot}[2]{%
	\begin{figure}[H]
		\centering
		\IfFileExists{#1}{%
			\includegraphics[width=0.92\textwidth]{#1}%
		}{%
			\fbox{\parbox[c][45mm][c]{0.9\textwidth}{\centering\itshape
				Место для скриншота\\[2pt]\ttfamily\small #1}}%
		}
		\caption{#2}
	\end{figure}}

\begin{document}

\begin{titlepage}
	\centering

	МОСКОВСКИЙ АВИАЦИОННЫЙ ИНСТИТУТ\par
	(НАЦИОНАЛЬНЫЙ ИССЛЕДОВАТЕЛЬСКИЙ УНИВЕРСИТЕТ)\par

	\vspace{6mm}

	Институт №8 <<Компьютерные науки и прикладная математика>>\par
	Кафедра 806 <<Вычислительная математика и программирование>>\par

	\vspace{6cm}

	{\Large Лабораторная работа №1\par}
	\vspace{3mm}
	{\large по курсу <<Компьютерная графика>>\par}
	\vspace{8mm}
	{\Large\bfseries Основы 3D графики\par}
	\vspace{3mm}
	{\large Вариант: правильный октаэдр\par}

	\vspace{4cm}

	\begin{flushright}
		\begin{minipage}{0.5\textwidth}
			Выполнил: \Student\par
			Группа: \Group\par
			\vspace{6mm}
			Преподаватель: \Teacher\par
		\end{minipage}
	\end{flushright}

	\vfill
	Москва, \Year\par
\end{titlepage}

\section{Условие}

Цель работы — познакомиться с основами трёхмерной графики: построением простых
3D-объектов, проекцией трёхмерной сцены на двумерную плоскость экрана, а также
научиться работать с матрицами перспективной и ортографической проекций и
матрицами аффинных преобразований.

Вариант задания — построить и вывести на экран правильный октаэдр.
Программа должна работать в реальном времени с использованием Vulkan, GLFW и
библиотеки интерфейса ImGui и допускать динамическую смену проекции и
преобразований объекта.

Выполнены 5 дополнительных заданий:

\begin{enumerate}[nosep, leftmargin=*]
	\item переключение матрицы проекции (ортографическая и перспективная) в интерфейсе;
	\item элементы интерфейса для изменения позиции, поворота и растяжения фигуры;
	\item движение и поворот фигуры по сложной траектории, пауза и воспроизведение
	анимации, изменение скорости и параметров траектории;
	\item изменение цвета фигуры с помощью элемента \texttt{ColorEdit};
	\item процедурное назначение цветов вершинам в зависимости от их позиции в
	локальной системе координат; выбранный в интерфейсе цвет умножается на цвет
	вершины.
\end{enumerate}


\section{Метод решения}

\subsection{Однородные координаты и аффинные преобразования}

Объект описывается один раз в собственной (локальной) системе координат, а
затем должен быть размещён в сцене, снят виртуальной камерой и спроецирован на
плоскость экрана. Каждое из этих действий — преобразование координат вершины.

Повороты и масштабирование линейны и выражаются матрицей $3\times3$. Перенос
линейным преобразованием не является: сдвиг на вектор $\mathbf{t}$ нельзя
записать как умножение вектора на матрицу $3\times3$. Поэтому точка
записывается в однородных координатах — четырёхмерным вектором
$(x, y, z, 1)^{\mathsf{T}}$, а преобразования — матрицами $4\times4$:

\begin{equation}
	T(\mathbf{t}) =
	\begin{pmatrix}
		1 & 0 & 0 & t_x \\
		0 & 1 & 0 & t_y \\
		0 & 0 & 1 & t_z \\
		0 & 0 & 0 & 1
	\end{pmatrix},
	\quad
	S(\mathbf{s}) =
	\begin{pmatrix}
		s_x & 0 & 0 & 0 \\
		0 & s_y & 0 & 0 \\
		0 & 0 & s_z & 0 \\
		0 & 0 & 0 & 1
	\end{pmatrix},
	\quad
	R_y(\theta) =
	\begin{pmatrix}
		\cos\theta & 0 & \sin\theta & 0 \\
		0 & 1 & 0 & 0 \\
		-\sin\theta & 0 & \cos\theta & 0 \\
		0 & 0 & 0 & 1
	\end{pmatrix}.
\end{equation}

Четвёртая координата решает две задачи. Во-первых, столбец переноса умножается
на единицу и просто прибавляется к результату. Во-вторых, после всех
преобразований выполняется перспективное деление: аппаратура вычисляет
$(x/w,\; y/w,\; z/w)$ — именно деление на $w$ даёт эффект перспективы.

Благодаря однородным координатам все преобразования становятся объектами
одного типа и перемножаются между собой, поэтому цепочка сворачивается в одну
матрицу. Она вычисляется один раз за кадр, а к каждой вершине применяется
однократно:

\begin{equation}
	\mathbf{v}' = P \cdot V \cdot M \cdot \mathbf{v},
	\label{eq:mvp}
\end{equation}

где $M$~— матрица модели, $V$~— матрица вида, $P$~— матрица проекции. Матрица
модели собирается как

\begin{equation}
	M = T(\mathbf{t}) \cdot R_z(\gamma) R_x(\alpha) R_y(\beta) \cdot S(\mathbf{s}).
	\label{eq:model}
\end{equation}

Порядок существенен: матрицы перемножаются справа налево, поэтому к вершине
сначала применяется масштабирование, затем поворот и только потом перенос. При
обратном порядке перенос был бы выполнен до поворота и фигура вращалась бы
вокруг начала координат сцены, а не вокруг собственного центра.

Камеры как отдельной сущности в математике преобразований нет: вместо
перемещения наблюдателя выполняется обратное перемещение всей сцены. Матрица
вида $V$ строится по положению камеры $\mathbf{e}$, точке наблюдения
$\mathbf{c}$ и вектору <<вверх>> $\mathbf{u}_0$ через ортонормированный базис
камеры:

\begin{equation}
	\mathbf{f} = \frac{\mathbf{c} - \mathbf{e}}{\|\mathbf{c} - \mathbf{e}\|}, \quad
	\mathbf{r} = \frac{\mathbf{f} \times \mathbf{u}_0}{\|\mathbf{f} \times \mathbf{u}_0\|}, \quad
	\mathbf{u} = \mathbf{r} \times \mathbf{f}.
\end{equation}

\subsection{Матрицы проекций}

\textbf{Перспективная проекция} моделирует зрение человека: чем дальше объект,
тем меньше он выглядит. Матрица строится по вертикальному углу обзора
$\varphi$, отношению сторон окна $a$ и расстояниям до ближней и дальней
плоскостей отсечения $n$ и $f$. Обозначив $k = \operatorname{ctg}(\varphi/2)$:

\begin{equation}
	P_{\text{persp}} =
	\begin{pmatrix}
		\dfrac{k}{a} & 0 & 0 & 0 \\[8pt]
		0 & -k & 0 & 0 \\[8pt]
		0 & 0 & \dfrac{f}{n - f} & \dfrac{n f}{n - f} \\[8pt]
		0 & 0 & -1 & 0
	\end{pmatrix}.
	\label{eq:persp}
\end{equation}

Ключевой элемент — $-1$ в последней строке. Он записывает в $w$ величину $-z$,
то есть расстояние от камеры до точки. При последующем делении на $w$ экранные
координаты уменьшаются обратно пропорционально расстоянию, что и даёт
перспективное сокращение.

\textbf{Ортографическая проекция} переносит точки на плоскость параллельно,
без схождения линий: размер объекта не зависит от расстояния до камеры. Если
$h$~— высота видимой области в мировых единицах, $t = h/2$ и $r = t a$, то

\begin{equation}
	P_{\text{ortho}} =
	\begin{pmatrix}
		\dfrac{1}{r} & 0 & 0 & 0 \\[8pt]
		0 & -\dfrac{1}{t} & 0 & 0 \\[8pt]
		0 & 0 & \dfrac{1}{n - f} & \dfrac{n}{n - f} \\[8pt]
		0 & 0 & 0 & 1
	\end{pmatrix}.
	\label{eq:ortho}
\end{equation}

Последняя строка здесь единичная: $w$ остаётся равным единице, деления не
происходит, и зависимость размера от расстояния исчезает. Таким образом,
различие двух проекций сводится к одному элементу матрицы, и задание~1
реализовано выбором функции, строящей $P$.

Обе формулы отличаются от приводимых в литературе по OpenGL по двум причинам.
Во-первых, в Vulkan ось $Y$ в координатах кадра направлена вниз, тогда как в
мировой системе координат она направлена вверх; поэтому обе матрицы содержат
знак минус во второй строке, иначе изображение оказалось бы перевёрнутым.
Во-вторых, канонический объём видимости в Vulkan ограничен по глубине
диапазоном $z \in [0, 1]$, а не $[-1, 1]$, как в OpenGL; отсюда иной вид
третьей строки.

Все матрицы реализованы самостоятельно в файле \texttt{math.hpp}, без
использования сторонних математических библиотек. Хранение — по столбцам, в том
же порядке, в котором язык GLSL интерпретирует тип \texttt{mat4}.

\subsection{Геометрия октаэдра}

Правильный октаэдр — правильный многогранник с шестью вершинами, двенадцатью
рёбрами и восемью гранями, каждая из которых является равносторонним
треугольником. Наиболее удобное представление — вершины на осях координат:
$(\pm 1, 0, 0)$, $(0, \pm 1, 0)$, $(0, 0, \pm 1)$. При таком выборе все рёбра
равны $\sqrt{2}$, а центр фигуры совпадает с началом координат, что удобно для
поворотов.

\begin{figure}[H]
	\centering
	\begin{tikzpicture}[scale=1.2, line join=round]
		\coordinate (v2) at (0, 2.2);
		\coordinate (v3) at (0, -2.2);
		\coordinate (v0) at (2, 0);
		\coordinate (v1) at (-2, 0);
		\coordinate (v4) at (0, -0.75);
		\coordinate (v5) at (0, 0.75);

		% невидимые рёбра, сходящиеся в дальней вершине
		\draw[dashed, gray] (v5) -- (v0);
		\draw[dashed, gray] (v5) -- (v1);
		\draw[dashed, gray] (v5) -- (v2);
		\draw[dashed, gray] (v5) -- (v3);

		\draw[thick] (v2) -- (v0) -- (v3) -- (v1) -- cycle;
		\draw[thick] (v2) -- (v4) -- (v3);
		\draw[thick] (v0) -- (v4) -- (v1);

		\foreach \p in {v0,v1,v2,v3,v4,v5}
			\fill (\p) circle (1.4pt);

		\node[right] at (v0) {$0\;(1,0,0)$};
		\node[left] at (v1) {$(-1,0,0)\;1$};
		\node[above] at (v2) {$2\;(0,1,0)$};
		\node[below] at (v3) {$3\;(0,-1,0)$};
		\node[below right=0pt and 2pt] at (v4) {$4\;(0,0,1)$};
		\node[above right=0pt and 2pt] at (v5) {$5\;(0,0,-1)$};
	\end{tikzpicture}
	\caption{Вершины октаэдра и их номера в вершинном буфере. Пунктиром
	показаны рёбра, сходящиеся в дальней вершине}
	\label{fig:octahedron}
\end{figure}

Грани задаются индексным буфером: восемь треугольников, по четыре при верхней
и нижней вершинах (таблица~\ref{tab:faces}). Использование индексного буфера
позволяет хранить шесть вершин вместо двадцати четырёх: каждая вершина
принадлежит четырём граням и описывается один раз.

\begin{table}[H]
	\caption{Грани октаэдра и направления их внешних нормалей}
	\label{tab:faces}
	\centering
	\vspace{2mm}
	\begin{tabular}{llll}
		\toprule
		Индексы & Нормаль & Индексы & Нормаль \\
		\midrule
		2, 4, 0 & $(+1, +1, +1)$ & 3, 0, 4 & $(+1, -1, +1)$ \\
		2, 1, 4 & $(-1, +1, +1)$ & 3, 4, 1 & $(-1, -1, +1)$ \\
		2, 5, 1 & $(-1, +1, -1)$ & 3, 1, 5 & $(-1, -1, -1)$ \\
		2, 0, 5 & $(+1, +1, -1)$ & 3, 5, 0 & $(+1, -1, -1)$ \\
		\bottomrule
	\end{tabular}
\end{table}

Для правильного изображения выпуклого тела используются два независимых
механизма. \emph{Буфер глубины}: вместе с цветом для каждого пикселя хранится
глубина ближайшего записанного фрагмента, а фрагмент, оказавшийся дальше,
отбрасывается. \emph{Отсечение нелицевых граней}: грань, обращённая от
наблюдателя, не растеризуется вообще; лицевая сторона определяется по знаку
площади треугольника в координатах кадра, то есть по направлению обхода его
вершин.

Порядок вершин в каждой тройке выбран так, чтобы обход был против часовой
стрелки при взгляде снаружи. Это проверяется по знаку скалярного произведения
нормали грани на радиус-вектор её центра: для всех восьми граней произведение
положительно, то есть нормали направлены наружу. Инверсия оси $Y$ в матрице
проекции переводит мировую систему координат, где ось $Y$ направлена вверх, в
систему координат кадра, где она направлена вниз, поэтому направление обхода в
терминах Vulkan сохраняется, и в конвейере задано
\texttt{frontFace = VK\_FRONT\_FACE\_COUNTER\_CLOCKWISE}.

\subsection{Используемые объекты Vulkan}

Для достижения цели работы создаются следующие объекты.

\begin{itemize}[leftmargin=*]
	\item \texttt{VkShaderModule} — два модуля, вершинный и фрагментный шейдеры,
	загружаемые из скомпилированных файлов SPIR-V. После создания конвейера
	уничтожаются, так как больше не нужны.

	\item \texttt{VkPipelineLayout} — описывает интерфейс между приложением и
	шейдерами; в данной работе содержит единственный диапазон push-констант.

	\item \texttt{VkPipeline} — графический конвейер. Существенные настройки:
	формат вершины из двух величин \texttt{VK\_FORMAT\_R32G32B32\_SFLOAT}
	(позиция и цвет); примитив — список треугольников; отсечение нелицевых
	граней с обходом против часовой стрелки; включённые тест и запись глубины с
	функцией сравнения <<меньше>>. Область вывода и прямоугольник отсечения
	объявлены динамическими и задаются при записи команд, поэтому при изменении
	размера окна конвейер не требуется пересоздавать.

	\item \texttt{VkBuffer} — вершинный и индексный буферы, размещённые через
	распределитель памяти \texttt{VmaAllocator} в памяти, доступной для записи
	со стороны процессора. Для объёма в несколько сотен байт промежуточный
	буфер в памяти видеокарты не требуется.

	\item \texttt{VkRenderPass} и \texttt{VkFramebuffer} — проход отрисовки с
	цветовым вложением и вложением глубины; предоставлены стартовым кодом.
\end{itemize}

Матрица \eqref{eq:mvp} и цвет передаются в шейдер через push-константы —
небольшой блок данных, записываемый непосредственно в командный буфер. Его
размер составляет 80 байт (матрица $4\times4$ из чисел одинарной точности и
четырёхкомпонентный цвет), что укладывается в гарантируемые спецификацией
128 байт. Такой способ не требует ни буферов равномерных данных, ни пулов
дескрипторов.

\subsection{Интерфейс, анимация и цвет}

Окно интерфейса разделено на четыре раздела, соответствующих дополнительным
заданиям: выбор проекции, преобразования фигуры, анимация и цвет. Значения
ползунков разделов \textbf{Transform} и \textbf{Animation} используются для
построения матрицы модели \eqref{eq:model}: таким образом, задания~2 и~3
сводятся к разным способам собрать $M$.

Траектория движения (задание~3) — вертикальная восьмёрка, намотанная на
окружность:

\begin{equation}
	x(t) = R\cos t, \qquad
	y(t) = H\sin 2t, \qquad
	z(t) = R\sin t,
\end{equation}

где $R$~— радиус окружности, $H$~— амплитуда вертикальных колебаний. Удвоенная
частота по оси $Y$ даёт два подъёма за один оборот, то есть траектория не
является плоской кривой. Параметр $t$ наращивается на величину
$\Delta t \cdot k$, где $\Delta t$~— длительность кадра, а $k$~— задаваемая в
интерфейсе скорость; при паузе наращивание прекращается. Независимо от движения
по траектории фигура вращается вокруг собственной оси с задаваемой угловой
скоростью.

Цвет вершины (задание~5) вычисляется из её позиции в локальной системе
координат при создании вершинного буфера: координата из диапазона $[-1, 1]$
переводится в канал цвета из диапазона $[0, 1]$ по формуле
$c = 0{,}5\,p + 0{,}5$. Выбранный в интерфейсе цвет (задание~4) умножается на
цвет вершины в вершинном шейдере, а интерполяция между вершинами выполняется
растеризатором. Флажок \texttt{Procedural vertex colors} заменяет цвета вершин
белым, и тогда фигура окрашивается только выбранным цветом.

\lstinputlisting[language=GLSL, caption={Вершинный шейдер: преобразование
вершины и расчёт цвета}]{shaders/octahedron.vert}

% ------------------------------------------------------------------
\section{Результаты}

\screenshot{images/perspective.png}{Октаэдр в перспективной проекции. Цвета
вершин заданы процедурно от их позиции в локальной системе координат и
интерполируются вдоль граней}

\screenshot{images/orthographic.png}{Та же сцена в ортографической проекции:
параллельные рёбра не сходятся, размер фигуры не зависит от удаления камеры}

\screenshot{images/transform.png}{Изменение позиции, поворота и растяжения
фигуры ползунками раздела \textbf{Transform}}

\screenshot{images/color.png}{Изменение цвета фигуры: выбранный в
\texttt{ColorEdit} цвет умножается на цвета вершин}

Различие проекций наиболее заметно, если сместить фигуру от центра экрана и
увеличить её: в перспективной проекции грани, расположенные дальше от камеры,
сокращаются, а параллельные рёбра сходятся; в ортографической проекции
параллельность рёбер сохраняется, и размер фигуры не зависит от удаления
камеры.

% ------------------------------------------------------------------
\section{Выводы}

В ходе работы я научился строить простые трёхмерные объекты и выводить их на
экран средствами Vulkan в реальном времени, работать с матрицами аффинных
преобразований и проекций, а также пользоваться буфером глубины и отсечением
нелицевых граней для правильного изображения выпуклого тела.

Главное, что я понял: представление точек в однородных координатах позволяет
описать перенос, поворот, масштабирование, видовое преобразование и проекцию
объектами одного типа — матрицами $4\times4$. Благодаря этому вся
последовательность преобразований сворачивается в одну матрицу, вычисляемую
один раз за кадр, а каждая вершина обрабатывается единственным умножением.
Практическим следствием стало то, что каждое из дополнительных заданий свелось
к изменению одной из матриц, а код отрисовки, шейдеры и буферы остались
неизменными: переключение проекции — это выбор функции, строящей матрицу $P$, а
преобразования и анимация — разные способы собрать матрицу $M$.

Отдельно отмечу особенности соглашений Vulkan, которые пришлось учесть:
диапазон глубины $[0, 1]$ вместо $[-1, 1]$ и направленная вниз ось $Y$ в
координатах кадра. Обе матрицы проекции содержат поправку на это; формулы,
взятые из литературы по OpenGL без такой поправки, дали бы перевёрнутое
изображение и неверную работу буфера глубины.

Полезным оказался и опыт отладки отсечения граней: ошибка в выборе лицевой
стороны не бросается в глаза, поскольку выпуклая фигура остаётся похожей на
себя — вместо ближних граней просто отображаются дальние. Обнаружить это
удалось только по цветам вершин, так как в центре фигуры оказался цвет дальней
вершины вместо ближней.

\end{document}
